% ============================================================
% DOCUMENTCLASS
% ============================================================
\documentclass[12pt,oneside]{book}

% ============================================================
% METADATOS
% ============================================================
\newcommand{\universidad}{Universidad de Buenos Aires (UBA)}
\newcommand{\facultad}{Facultad de Derecho}
\newcommand{\carrera}{Orientación Notarial y Registral}
\newcommand{\materia}{Derechos Reales}
\newcommand{\titulo}{Derecho Real de Servidumbre}
\newcommand{\autor}{Isaac Edgar Camacho Ocampo}
\newcommand{\anio}{2026}

% ============================================================
% PAQUETES
% ============================================================
\usepackage[utf8]{inputenc}
\usepackage[T1]{fontenc}
\usepackage[spanish,es-nodecimaldot]{babel}

\usepackage[
    top=2.5cm,
    bottom=2.5cm,
    left=3cm,
    right=2.5cm,
    headheight=15pt
]{geometry}

\usepackage{amsmath}
\usepackage{amssymb}
\usepackage{mathtools}

\usepackage{graphicx}
\usepackage{float}
\usepackage{caption}
\usepackage{subcaption}

\usepackage{booktabs}
\usepackage{array}
\usepackage{longtable}
\usepackage{tabularx}

\usepackage{enumitem}

\usepackage{xcolor}
\usepackage[most]{tcolorbox}

\usepackage{fancyhdr}
\usepackage{ragged2e}

\usepackage[
    colorlinks=true,
    linkcolor=blue!50!black,
    citecolor=green!40!black,
    urlcolor=blue!60!black,
    pdftitle={\titulo},
    pdfauthor={\autor},
    pdfsubject={Apuntes universitarios},
    bookmarks=true
]{hyperref}

% ============================================================
% CONFIGURACIÓN
% ============================================================
\graphicspath{
    {./img/}
    {./graficos/}
    {./figuras/}
}

\setcounter{tocdepth}{2}

\setlength{\parindent}{0pt}
\setlength{\parskip}{0.5em}

\setlist[itemize]{
    topsep=0.3em,
    itemsep=0.2em
}

\setlist[enumerate]{
    topsep=0.3em,
    itemsep=0.2em
}

\clubpenalty=10000
\widowpenalty=10000

% ============================================================
% COMANDOS
% ============================================================
\newcommand{\abs}[1]{\left|#1\right|}
\newcommand{\norm}[1]{\left\lVert#1\right\rVert}
\newcommand{\vect}[1]{\mathbf{#1}}
\newcommand{\deriv}[2]{\frac{d#1}{d#2}}
\newcommand{\pderiv}[2]{\frac{\partial#1}{\partial#2}}

% ============================================================
% ENTORNOS / CAJAS
% ============================================================
\newtcolorbox{definition}[1]{
    enhanced,
    breakable,
    title={Definición: #1},
    fonttitle=\bfseries,
    colback=blue!3,
    colframe=blue!50!black,
    boxrule=0.8pt,
    arc=2mm
}

\newtcolorbox{important}{
    enhanced,
    breakable,
    title={Importante},
    fonttitle=\bfseries,
    colback=yellow!8,
    colframe=orange!70!black,
    boxrule=0.8pt,
    arc=2mm
}

\newtcolorbox{example}[1]{
    enhanced,
    breakable,
    title={Ejemplo: #1},
    fonttitle=\bfseries,
    colback=green!4,
    colframe=green!40!black,
    boxrule=0.8pt,
    arc=2mm
}

\newtcolorbox{formula}[1]{
    enhanced,
    breakable,
    title={Fórmula / Regla: #1},
    fonttitle=\bfseries,
    colback=purple!4,
    colframe=purple!50!black,
    boxrule=0.8pt,
    arc=2mm
}

\newtcolorbox{exercise}[1]{
    enhanced,
    breakable,
    title={Ejercicio: #1},
    fonttitle=\bfseries,
    colback=gray!5,
    colframe=gray!60!black,
    boxrule=0.8pt,
    arc=2mm
}

\newtcolorbox{note}{
    enhanced,
    breakable,
    title={Nota},
    fonttitle=\bfseries,
    colback=white,
    colframe=gray!50,
    boxrule=0.6pt,
    arc=2mm
}

% ============================================================
% CONFIGURACIÓN FINAL
% ============================================================
\pagestyle{fancy}
\fancyhf{}
\fancyhead[L]{\nouppercase{\leftmark}}
\fancyhead[R]{\materia}
\fancyfoot[C]{\thepage}
\renewcommand{\headrulewidth}{0.4pt}
\renewcommand{\footrulewidth}{0pt}

\fancypagestyle{plain}{
    \fancyhf{}
    \fancyfoot[C]{\thepage}
    \renewcommand{\headrulewidth}{0pt}
}

\renewcommand{\arraystretch}{1.25}

% ============================================================
% DOCUMENT
% ============================================================
\begin{document}

% ------------------------------------------------------------
% PORTADA
% ------------------------------------------------------------
\begin{titlepage}

\thispagestyle{empty}

\begin{center}

\vspace*{1cm}

{\large \textsc{\universidad}}

\vspace{0.2cm}

{\large \textsc{\facultad}}

\vspace{0.2cm}

{\large \carrera}

\vspace{3cm}

{\Huge \textbf{\materia}}

\vspace{1cm}

{\LARGE \titulo}

\vspace{0.8cm}

{\small \textit{Comprender es el primer paso para convertir el estudio en conocimiento.}}

\vspace{1.2cm}

\rule{0.70\textwidth}{0.5pt}

\vspace{0.5cm}

{\large Apuntes de estudio}

\vspace{0.5cm}

\rule{0.70\textwidth}{0.5pt}

\vfill

{\large \autor}

\vspace{0.3cm}

{\large \anio}

\end{center}

\vfill

\begin{flushleft}
\textit{Este es un modesto aporte a los estudiantes de ``Derechos Reales'' no pretende sustituir el material oficial de la materia.}
\end{flushleft}

\end{titlepage}

% ------------------------------------------------------------
% ÍNDICE
% ------------------------------------------------------------
\tableofcontents

% ------------------------------------------------------------
% CONTENIDO
% ------------------------------------------------------------

% ============================================================
\chapter[Concepto y objeto]{Concepto y objeto de la servidumbre}
% ============================================================

\section{Concepto}

La servidumbre es un \textbf{derecho real de disfrute sobre cosa ajena} que \textbf{vincula dos fundos entre sí}.

A partir de la constitución del derecho de servidumbre, el dominio del titular de uno de los fundos se \textbf{imperfecciona}: el dominio, que era absoluto, pasa a ser un \textbf{dominio imperfecto}, porque la constitución de la servidumbre \textbf{desmembra el dominio} del titular.

\begin{note}
Para presentar la figura se utilizó un esquema básico, a modo de pizarra, en el que las tierras y los diversos fundos (campos o potreros) figuran en naranja. Se emplea este esquema como recurso de presentación y no porque los fundos deban ser necesariamente agrícolas. En el esquema se identifican el Fundo A y el Fundo B, una parte del terreno delimitada con líneas rojas y la Ruta Nacional N.º 2 como vía pública. La imagen del esquema no se encuentra disponible en las notas originales; se conserva su descripción textual.
% TODO: imagen del esquema (Fundo A, Fundo B, Ruta Nacional N.º 2) no disponible en las notas originales
\end{note}

\section{Los dos fundos: dominante y sirviente}

Se trata de un derecho real que vincula a dos fundos entre sí. En el esquema, el \textbf{Fundo A} adquiere la calidad de \textbf{fundo dominante} y el \textbf{Fundo B} la calidad de \textbf{fundo sirviente}.

Ambos fundos se vinculan a partir de la constitución del derecho real, que implica que el titular del fundo sirviente (Fundo B) deberá prestar a favor del otro fundo cierta \textbf{utilidad}. Esa utilidad puede consistir en pasar, transitar o sacar agua, e incluso puede ser de mero recreo.

En todos los casos, la utilidad se traduce en la obligación de dejar que el titular del fundo dominante haga algo, o de abstenerse de realizar determinada acción o actividad que el titular del dominio tendría si no existiese la servidumbre.

\begin{important}
\textbf{En ningún caso puede tratarse de una obligación de hacer.} La carga siempre consiste en una abstención o en permitir que el otro haga.
\end{important}

\section{Definición legal}

\begin{definition}{Servidumbre (Art. 2162 CCCN)}
La servidumbre es el derecho real que se establece entre dos inmuebles y que concede al titular del inmueble dominante determinada utilidad sobre el inmueble del fondo sirviente, que es el inmueble ajeno.
\end{definition}

\section{Objeto de la servidumbre}

El objeto de la servidumbre puede ser todo el terreno o una parte determinada del terreno (por ejemplo, si se tratara de transitar, la franja identificada en el esquema con las líneas rojas).

Siempre se trata de \textbf{dos inmuebles con distinta titularidad}, que tienen por objeto una parte material determinada o todo el inmueble.

\section{Cuadro Cornell}

\begin{table}[H]
\centering
\begin{tabularx}{\textwidth}{
    !{\vrule width 0.4pt}>{\raggedright\arraybackslash}p{0.26\textwidth}
    !{\vrule width 0.4pt}>{\raggedright\arraybackslash}X
    !{\vrule width 0.4pt}
}
\toprule
\textbf{Preguntas / palabras clave} &
\textbf{Notas / respuestas} \\
\midrule

\textbf{¿Qué es la servidumbre y qué establece el Art. 2162 CCCN?} &
Es un derecho real de disfrute sobre cosa ajena que vincula dos fundos entre sí. Según el Art. 2162, se establece entre dos inmuebles y concede al titular del inmueble dominante determinada utilidad sobre el inmueble del fondo sirviente, que es el inmueble ajeno. \\

\textbf{¿Qué diferencia hay entre fundo dominante y fundo sirviente?} &
El fundo dominante es el que recibe la utilidad; el fundo sirviente es el que debe prestarla, y su titular soporta la carga. Siempre se trata de dos inmuebles con distinta titularidad. \\

\textbf{¿Qué utilidad puede conceder la servidumbre y en qué se traduce para el fundo sirviente?} &
La utilidad puede ser pasar, transitar, sacar agua o incluso de mero recreo. Para el titular del fundo sirviente se traduce en dejar que el titular del dominante haga, o en abstenerse de una acción que tendría si no existiese la servidumbre. Nunca es una obligación de hacer. \\

\textbf{¿Por qué la servidumbre imperfecciona el dominio?} &
Porque su constitución desmembra el dominio del titular del fundo: el dominio, que era absoluto, pasa a ser un dominio imperfecto. \\

\textbf{¿Sobre qué recae la servidumbre?} &
Sobre todo el inmueble o sobre una parte material determinada de él. \\

\midrule

\multicolumn{2}{!{\vrule width 0.4pt}>{\raggedright\arraybackslash}X!{\vrule width 0.4pt}}{%
\textbf{Resumen:} La servidumbre vincula dos inmuebles de distinta titularidad: el dominante recibe una utilidad y el sirviente la presta, siempre mediante una abstención o un permitir (nunca un hacer). Su constitución desmembra el dominio del fundo sirviente, que se vuelve imperfecto, y puede recaer sobre todo el inmueble o sobre una parte determinada.
} \\

\bottomrule
\end{tabularx}
\end{table}

% ============================================================
\chapter{Clasificación de las servidumbres}
% ============================================================

\section{Servidumbres positivas y negativas}

La denominación que se da al tipo de utilidad permite clasificar las servidumbres en positivas y negativas.

\begin{table}[H]
\centering
\begin{tabularx}{\textwidth}{>{\raggedright\arraybackslash}X >{\raggedright\arraybackslash}X}
\toprule
\textbf{Positivas} & \textbf{Negativas} \\
\midrule
La carga consiste en \textbf{soportar} el ejercicio de una acción.
&
El titular del fundo sirviente debe \textbf{abstenerse} de realizar alguna facultad. \\[0.8em]
\textit{Ejemplo:} soportar que pasen por todo el fundo o por un sector determinado de él.
&
\textit{Ejemplo:} servidumbre de vista; el titular sirviente no puede construir y afectar la vista del fundo dominante. \\
\bottomrule
\end{tabularx}
\end{table}

\section{Servidumbres reales y personales}

Esta clasificación \textbf{no se refiere a la distinción entre derechos reales y derechos personales}, que son dos esferas diferentes de derechos. Las servidumbres son \textbf{siempre reales}: no existen servidumbres personales como instituto dentro de los derechos reales. Se trata de un instituto propio de los derechos reales que vincula dos fundos.

Se denomina \textbf{personales} a aquellas servidumbres (siempre reales) constituidas, ya sea por voluntad de las partes o por ley, a favor de una persona determinada, teniendo en cuenta quién es el titular del fundo dominante.
% TODO: contenido ambiguo en las notas originales: se afirma que la servidumbre personal puede constituirse «por voluntad de las partes o por ley», pero más adelante se la describe como siempre convencional.

En cambio, la \textbf{servidumbre real} se constituye teniendo en cuenta el fundo, con independencia de quién sea su titular.

\section{Servidumbres forzosas (de origen legal)}

En el caso de la servidumbre real puede tratarse de fundos encerrados que no tienen acceso a la vía pública, o de fundos que carecen de agua. En estos casos la servidumbre es de \textbf{origen legal}, porque el orden público interviene rigurosamente y exige e impone la servidumbre.

El orden público no admite, por resultar desfavorable para todos y para la comunidad, un fundo que carezca de toda utilidad. El ejemplo concreto es la servidumbre de paso o de tránsito.

\begin{example}{Fundo encerrado (servidumbre legal de tránsito)}
Se tienen dos fundos, y todo lo marcado en naranja es tierra. Uno de ellos no tiene salida a la vía pública (en el ejemplo, la Ruta Nacional N.º 2); no hay días de circulación. Por lo tanto, sus titulares no podrían ingresar y, si lograran hacerlo por algún medio aéreo, tampoco podrían salir. No tendría sentido. Lo mismo ocurre con los fundos que carecen de agua.

La servidumbre de origen legal impone necesariamente al titular del fundo que tiene una salida más directa a la vía pública (en este caso, a la Ruta Nacional N.º 2) tolerar y permitir que el titular del fundo dominante pase.
\end{example}

\begin{important}
La mirada del orden público está puesta en la \textbf{utilidad del fundo}, no en la calidad personal del titular de dominio. Cualquiera sea el titular, aun cuando mute el titular registral por enajenación o por transmisión \textit{mortis causa} a herederos, tendrá que poder pasar.
\end{important}

\section{Servidumbres personales convencionales}

Se analiza ahora una servidumbre de tránsito también, pero personal. En el esquema, el inmueble A tiene salida a la vía pública (la Ruta Nacional N.º 2) por un camino lateral, un camino rural que bordea los campos y continúa; recién a muchos kilómetros sale a la Ruta Nacional N.º 2.

No puede imponerse al titular del otro fundo que permita que A pase por su fundo solo porque le resulte más cómodo o el acceso sea más directo. El orden público no está vinculado aquí: solo impone la servidumbre cuando \textbf{no existe salida a la vía pública o no existe acceso al agua}.

\begin{example}{Servidumbre personal convencional}
El titular de A nunca tuvo dificultades y le resultó suficiente salir por el camino existente, con su caballo o su vehículo, hacia la ruta. Un día decide cambiar su actividad o vende, y el adquirente quiere realizar una actividad que impone un gran tránsito de camiones; ese camino le resulta poco práctico y además debe recorrer muchos kilómetros para llegar a la vía pública.

Entonces, el adquirente pacta con el titular del fundo sirviente la constitución de una servidumbre. Hay libertad y consenso: le propone un camino determinado por el cual pueda pasar con sus camiones. No hay imposición legal; es exclusivamente convencional. Las partes pactarán el precio y la duración.

Si no pudieran ponerse de acuerdo en la indemnización, A podría solicitarla judicialmente y el precio sería fijado por el juez. Nada de esto ocurre con la servidumbre personal cuando el orden público está satisfecho: la salida a la vía pública existe y la servidumbre es del interés de A, y del interés de B si le cierra la ecuación económica.
\end{example}
% TODO: contenido ambiguo en las notas originales: se indica que A podría solicitar judicialmente la indemnización y que el juez fijaría el precio en una servidumbre convencional, y luego que la servidumbre personal nunca puede constituirse judicialmente.

\section{Cuadro comparativo: servidumbre real y personal}

\begin{table}[H]
\centering
\small
\begin{tabularx}{\textwidth}{>{\raggedright\arraybackslash}p{0.2\textwidth} >{\raggedright\arraybackslash}X >{\raggedright\arraybackslash}X}
\toprule
\textbf{Aspecto} & \textbf{Real (forzosa)} & \textbf{Personal (convencional)} \\
\midrule
\textbf{Origen} & Legal; orden público & Convencional; voluntad de las partes \\
\textbf{Requisito} & Fundo encerrado o sin agua & Solo el interés de las partes \\
\textbf{Camino} & El más directo a la vía pública & El que acuerden las partes \\
\textbf{Precio} & Fijado por el juez si no hay acuerdo & Libremente pactado \\
\textbf{Transmisión} & Con el inmueble (muta el titular) & Se extingue al transmitir \\
\textbf{Duración} & Mientras subsiste la utilidad & Temporal; la acordada \\
\textbf{Plazo máximo sin estipular} & Indeterminado (mientras sea útil) & Vitalicia (persona humana); 50 años (persona jurídica) \\
\bottomrule
\end{tabularx}
\end{table}

\section{Cuadro Cornell}

\begin{table}[H]
\centering
\begin{tabularx}{\textwidth}{
    !{\vrule width 0.4pt}>{\raggedright\arraybackslash}p{0.26\textwidth}
    !{\vrule width 0.4pt}>{\raggedright\arraybackslash}X
    !{\vrule width 0.4pt}
}
\toprule
\textbf{Preguntas / palabras clave} &
\textbf{Notas / respuestas} \\
\midrule

\textbf{¿Cuándo es positiva y cuándo negativa la servidumbre?} &
Es positiva cuando la carga consiste en soportar el ejercicio de una acción (por ejemplo, que pasen por el fundo). Es negativa cuando el titular del fundo sirviente debe abstenerse de realizar alguna facultad (por ejemplo, la servidumbre de vista: no puede construir y afectar la vista del dominante). \\

\textbf{¿Qué diferencia hay entre servidumbre real y personal?} &
La clasificación no se refiere a derechos reales versus personales: las servidumbres son siempre reales. La real se constituye teniendo en cuenta el fundo, con independencia de su titular, y se transmite con el inmueble. La personal se constituye teniendo en cuenta a la persona del titular del fundo dominante y se extingue al transmitirse el inmueble. \\

\textbf{¿Cuándo existe una servidumbre forzosa (legal) y por qué?} &
Cuando un fundo carece de salida a la vía pública (fundo encerrado) o de acceso al agua. El orden público no admite un fundo carente de toda utilidad, e impone la servidumbre al titular del fundo con salida más directa. La mirada está puesta en la utilidad del fundo y no en la calidad del titular. \\

\textbf{¿Puede imponerse una servidumbre solo porque resulta más cómoda?} &
No. El orden público solo interviene cuando no hay salida a la vía pública o acceso al agua. Si la salida existe, la servidumbre solo puede ser convencional: las partes acuerdan libremente camino, precio y duración. \\

\textbf{¿En qué se diferencian la servidumbre real (forzosa) y la personal (convencional)?} &
Origen (legal o convencional), requisito (fundo encerrado o sin agua, o interés de las partes), camino (el más directo o el acordado), precio (fijado por el juez si no hay acuerdo o libremente pactado), transmisión (con el inmueble o se extingue al transmitir), duración (mientras subsiste la utilidad o la acordada) y plazo máximo sin estipular (indeterminado o vitalicia/50 años). \\

\midrule

\multicolumn{2}{!{\vrule width 0.4pt}>{\raggedright\arraybackslash}X!{\vrule width 0.4pt}}{%
\textbf{Resumen:} Las servidumbres se clasifican en positivas (soportar) y negativas (abstenerse), y en reales (atienden al fundo y se transmiten con el inmueble) y personales (atienden al titular y se extinguen al transmitirse). Las forzosas o legales proceden por orden público cuando el fundo está encerrado o carece de agua; las convencionales nacen del acuerdo de partes sobre camino, precio y duración.
} \\

\bottomrule
\end{tabularx}
\end{table}

% ============================================================
\chapter[Constitución y legitimación]{Constitución y legitimación}
% ============================================================

\section{Forma: escritura pública}

La servidumbre siempre debe constituirse mediante \textbf{escritura pública}, porque se trata de un derecho real constituido sobre inmuebles. La escritura pública es el \textbf{título suficiente} en la constitución de la servidumbre.

\section{Duración de la servidumbre personal y presunción de vitalicidad}

Como la servidumbre personal es siempre convencional en su origen, tendrá una duración: es necesariamente convencional y temporal. Si no se determina el plazo en el instrumento constitutivo (la escritura pública), cuando se trata de persona humana se considera que es \textbf{vitalicia}.

\begin{quote}
\textbf{Art. 2165 CCCN.} \textit{Se presume que la servidumbre personal es vitalicia cuando no se ha determinado un plazo en el instrumento constitutivo. En el caso de persona jurídica, el plazo máximo es de cincuenta (50) años.}
\end{quote}

También en la servidumbre forzosa (la servidumbre legal: de tránsito, de fondos encerrados, de acueducto, de sacar agua), en todos los casos el titular del fundo sirviente tendrá el derecho de recibir la indemnización y no sufrir perjuicios.

\begin{important}
En ningún caso la servidumbre personal puede ser constituida judicialmente.
\end{important}

\section{Legitimación: quién puede constituir la servidumbre}

Puede constituir la servidumbre el \textbf{titular de dominio} y, en caso de condominio, \textbf{todos los condóminos} que totalicen el cien por ciento de la titularidad de la cosa.

\subsection{El nudo propietario y el usufructo}

Si el titular del fundo (B) hubiese constituido un usufructo, no puede a su vez constituir una servidumbre. El nudo propietario solo tiene el poder de disposición, y la servidumbre implica un menoscabo a ese derecho, que él mismo ha desplazado a favor del usufructuario.

Cuando existe un usufructo, no es el titular de dominio (el nudo propietario) quien tiene la facultad: el \textbf{usufructuario} sí podrá constituir la servidumbre, y esta estará vinculada con el tiempo que dure el usufructo, que es siempre un derecho temporal.

\subsection{Uso, habitación y anticresis}

Si el nudo propietario hubiese constituido, en lugar de un usufructo, un derecho de uso, de habitación o de anticresis, podrá constituir la servidumbre en cualquiera de estas dos formas:
\begin{itemize}
    \item con la autorización del habitador, del usuario o del acreedor anticresista; o
    \item constituyéndola de modo que comience a regir como condición necesaria cuando se hayan extinguido el uso, la habitación o el derecho de anticresis.
\end{itemize}

\section{Cuadro Cornell}

\begin{table}[H]
\centering
\begin{tabularx}{\textwidth}{
    !{\vrule width 0.4pt}>{\raggedright\arraybackslash}p{0.26\textwidth}
    !{\vrule width 0.4pt}>{\raggedright\arraybackslash}X
    !{\vrule width 0.4pt}
}
\toprule
\textbf{Preguntas / palabras clave} &
\textbf{Notas / respuestas} \\
\midrule

\textbf{¿Qué forma es necesaria para constituir una servidumbre?} &
Siempre escritura pública, por tratarse de un derecho real sobre inmuebles. Este instrumento es el título suficiente de constitución. \\

\textbf{¿Quién está legitimado para constituirla?} &
El titular de dominio, o todos los condóminos que totalicen el cien por ciento de la titularidad de la cosa. \\

\textbf{¿Puede el nudo propietario constituirla durante el usufructo?} &
No: solo tiene el poder de disposición y la servidumbre menoscaba el derecho que él desplazó a favor del usufructuario. Quien puede constituirla es el usufructuario, por el tiempo que dure el usufructo, que es siempre temporal. \\

\textbf{¿Qué ocurre si existe uso, habitación o anticresis?} &
La servidumbre puede constituirse con la autorización del habitador, usuario o acreedor anticresista, o bien comenzar a regir una vez extinguidos esos derechos. \\

\textbf{Art. 2165: ¿qué plazo tiene la servidumbre personal sin plazo determinado?} &
Se presume vitalicia si no se determinó plazo en el instrumento constitutivo; si se trata de persona jurídica, el plazo máximo es de cincuenta años. La servidumbre personal nunca puede ser constituida judicialmente. \\

\textbf{¿Qué derecho tiene el titular del fundo sirviente en las servidumbres forzosas?} &
Recibir la indemnización y no sufrir perjuicios. \\

\midrule

\multicolumn{2}{!{\vrule width 0.4pt}>{\raggedright\arraybackslash}X!{\vrule width 0.4pt}}{%
\textbf{Resumen:} La servidumbre se constituye siempre por escritura pública. Están legitimados el titular de dominio o todos los condóminos; el nudo propietario no puede constituirla durante el usufructo (sí el usufructuario, por el tiempo del usufructo), y ante uso, habitación o anticresis se requiere autorización o diferir sus efectos a la extinción de esos derechos. La servidumbre personal sin plazo es vitalicia (persona humana) o de hasta 50 años (persona jurídica) y nunca puede constituirse judicialmente.
} \\

\bottomrule
\end{tabularx}
\end{table}

% ============================================================
\chapter[Transmisión, derechos y obligaciones]{Transmisión, derechos y obligaciones}
% ============================================================

\section{Transmisión de las servidumbres reales}

Las servidumbres \textbf{reales}, que son aquellas que vinculan los fundos entre sí, se transmiten conjuntamente con el inmueble, ya enajene el inmueble el titular del fundo dominante o el titular del fundo sirviente. Como los vinculados son los fundos, quienquiera que sea el titular de uno u otro, la servidumbre se transmite con independencia de la mutación del titular.

\begin{quote}
\textbf{Art. 2172 CCCN.} \textit{Las servidumbres reales se transmiten conjuntamente con el inmueble al que acceden, tanto la constituida a su favor como la que lo grava. La servidumbre personal se extingue con la transmisión del inmueble.}
\end{quote}

Cuando se enajena el inmueble, los adquirentes lo reciben con la servidumbre. La servidumbre \textbf{personal}, en cambio, se extingue con la transmisión del inmueble, porque tanto en su constitución como en su duración se ha tenido en cuenta a esa persona determinada, que es el titular de A.

\section{Derechos y obligaciones del titular del fundo dominante}

\begin{quote}
\textbf{Arts. 2173 a 2175 CCCN.} \textit{El titular del fundo dominante puede realizar en el inmueble sirviente todas las obras necesarias para el ejercicio y conservación de la servidumbre. Tiene a su cargo las mejoras necesarias, como mantener el camino en condiciones. Debe comunicar al titular del fundo sirviente cualquier perturbación, de hecho o de derecho, de que tome conocimiento, para que pueda ejercer la defensa.}
\end{quote}

El titular del fundo dominante tiene la utilidad, pero debe realizar todas las mejoras necesarias y tiene la facultad de hacer en el inmueble sirviente todo aquello que sea necesario para que su derecho no se vea menoscabado:
\begin{itemize}
    \item si se trata de un camino, debe repararlo y mantenerlo en condiciones;
    \item si se hubiese producido un hecho natural, como la caída de árboles que le impidan salir, puede quitarlos;
    \item debe comunicar al titular del fundo sirviente cualquier perturbación, de hecho o de derecho, para que pueda ejercer la defensa.
\end{itemize}

\section{Inembargabilidad de la servidumbre}

\begin{quote}
\textbf{Art. 2178 CCCN.} \textit{La servidumbre no puede ser embargada ni vendida con independencia del inmueble al que accede, ya sea el inmueble dominante o el sirviente.}
\end{quote}

A diferencia de otros derechos reales ya estudiados, el Art. 2178 indica con claridad que en ningún caso la servidumbre en sí misma es ejecutable: no puede ser embargada ni vendida con independencia del fundo que grava o del fundo que beneficia.
% TODO: contenido ambiguo en las notas originales: se menciona como ejemplo de derecho ejecutable el caso «de los frutos», sin mayor precisión.

\section{Derechos y obligaciones del titular del fundo sirviente}

\begin{quote}
\textbf{Arts. 2180 y 2181 CCCN.} \textit{El titular del fundo sirviente conserva la disposición jurídica y material del inmueble. Puede ejercer todas las facultades propias de su derecho de dominio, con el límite de no afectar el ejercicio de la servidumbre. Puede realizar obras, constituir otros derechos reales y mejoras, siempre que no impidan o menoscaben el derecho del titular del fundo dominante.}
\end{quote}

Al titular del fundo sirviente, al que cuando la servidumbre es de origen legal o forzoso se le impone una carga como el acueducto o el tránsito, le quedan todas las facultades materiales y jurídicas propias de su derecho de dominio. El \textbf{límite} es que no puede realizar nada que afecte el derecho del titular del fundo dominante.

La afección o imperfección de su dominio desmembrado consiste en que no puede constituir otros derechos reales ni derechos personales que de alguna manera afecten el goce (la utilidad) que el titular del fundo dominante recibe de la servidumbre. Puede realizar obras, constituir otros derechos reales y hacer mejoras, pero en ningún caso hacer nada que afecte al titular del fundo dominante en el ejercicio de las facultades que le concede la servidumbre.

\section{Cuadro Cornell}

\begin{table}[H]
\centering
\begin{tabularx}{\textwidth}{
    !{\vrule width 0.4pt}>{\raggedright\arraybackslash}p{0.26\textwidth}
    !{\vrule width 0.4pt}>{\raggedright\arraybackslash}X
    !{\vrule width 0.4pt}
}
\toprule
\textbf{Preguntas / palabras clave} &
\textbf{Notas / respuestas} \\
\midrule

\textbf{Art. 2172: ¿cómo se transmiten las servidumbres reales y las personales?} &
Las reales se transmiten conjuntamente con el inmueble, tanto la constituida a su favor como la que lo grava, con independencia de la mutación del titular. La personal se extingue con la transmisión del inmueble, porque se tuvo en cuenta a la persona determinada titular de A. \\

\textbf{Arts. 2173 a 2175: ¿qué derechos y obligaciones tiene el titular del fundo dominante?} &
Puede realizar en el sirviente las obras necesarias para el ejercicio y conservación de la servidumbre. Tiene a su cargo las mejoras necesarias (por ejemplo, mantener el camino). Debe comunicar al sirviente cualquier perturbación, de hecho o de derecho, para que pueda ejercer la defensa. \\

\textbf{Art. 2178: ¿puede embargarse la servidumbre en forma autónoma?} &
No. No puede ser embargada ni vendida con independencia del inmueble al que accede, sea el dominante o el sirviente. \\

\textbf{Arts. 2180 y 2181: ¿qué conserva el titular del fundo sirviente y cuál es su límite?} &
Conserva la disposición jurídica y material del inmueble y todas las facultades propias del dominio. Puede realizar obras, constituir otros derechos reales y mejoras. Su límite es no afectar ni menoscabar el ejercicio de la servidumbre por parte del titular del fundo dominante. \\

\textbf{¿En qué consiste la imperfección del dominio del fundo sirviente?} &
En que no puede constituir otros derechos reales ni derechos personales que afecten el goce o utilidad que recibe el titular del fundo dominante. \\

\midrule

\multicolumn{2}{!{\vrule width 0.4pt}>{\raggedright\arraybackslash}X!{\vrule width 0.4pt}}{%
\textbf{Resumen:} Las servidumbres reales se transmiten con el inmueble; la personal se extingue al transmitirse. El titular dominante puede realizar las obras necesarias, tiene a su cargo las mejoras necesarias y debe comunicar las perturbaciones. El sirviente conserva las facultades del dominio con el límite de no afectar la servidumbre. La servidumbre es inembargable e inseparable del inmueble (Art. 2178).
} \\

\bottomrule
\end{tabularx}
\end{table}

% ============================================================
\chapter{Extinción de la servidumbre}
% ============================================================

\section{Causales de extinción}

Cuando se trata de fundos encerrados o sin agua, la servidumbre se mantiene en tanto se mantengan esas condiciones de encerramiento o la falta de agua.

\begin{quote}
\textbf{Art. 2182 CCCN (causales de extinción).} \textit{Son causales de extinción de la servidumbre: la desaparición de toda utilidad para el fundo dominante; la falta de uso por diez (10) años; en las personales, la muerte del titular cuando no se ha fijado un plazo (o, cuando sea persona jurídica, el vencimiento del plazo máximo de cincuenta (50) años); y el vencimiento del plazo fijado en el instrumento constitutivo.}
\end{quote}

\begin{itemize}
    \item \textbf{Desaparición de toda utilidad para el fundo dominante.} Si el fundo está encerrado y necesita salir a la vía pública, carecería de sentido mantener la servidumbre si luego se construyera un camino lateral: el interés del orden público desaparece y la servidumbre carece de objeto.
    \item \textbf{Falta de uso por diez años.}
    \item \textbf{En las servidumbres personales, la muerte del titular}, cuando no se haya fijado un plazo. La servidumbre personal es siempre temporal, y lo temporal por excelencia es la vida humana.
    \item \textbf{En el caso de persona jurídica}, el plazo máximo de cincuenta años, cuando no se ha fijado un plazo especial.
    \item \textbf{Vencimiento del plazo} fijado en el instrumento constitutivo.
\end{itemize}

\section{Efectos de la extinción}

Al extinguirse la servidumbre se extinguen todos los derechos personales o reales que pudieran estar vinculados con ella.

El titular del fundo sirviente, que veía aminorado su derecho, lo restablece: pasa a tener nuevamente un \textbf{dominio pleno} (también llamado \textbf{dominio perfecto}), porque ha desaparecido la carga que desmembraba su derecho.

\section{Síntesis de la unidad}

La servidumbre es un derecho real que siempre implica la vinculación entre dos fundos.

\begin{itemize}
    \item En las \textbf{servidumbres forzosas o legales} se tiene en cuenta la calidad del fundo con independencia del titular. Deben ser constituidas por imperio de la ley, más allá de la voluntad del titular del fundo sirviente, porque el orden público ha puesto allí su mirada.
    \item Las \textbf{servidumbres personales} son siempre reales, pero se las llama personales porque en su constitución se tiene en cuenta la persona del titular. Por eso son convencionales: las partes acuerdan el monto de la indemnización, la manera en que se pasará por el inmueble, y los derechos y obligaciones que surgen de una situación consensual que las vincula durante el tiempo que ambas hayan fijado.
\end{itemize}

\section{Cuadro Cornell}

\begin{table}[H]
\centering
\begin{tabularx}{\textwidth}{
    !{\vrule width 0.4pt}>{\raggedright\arraybackslash}p{0.26\textwidth}
    !{\vrule width 0.4pt}>{\raggedright\arraybackslash}X
    !{\vrule width 0.4pt}
}
\toprule
\textbf{Preguntas / palabras clave} &
\textbf{Notas / respuestas} \\
\midrule

\textbf{Art. 2182: ¿cuáles son las causales de extinción?} &
(1) Desaparición de toda utilidad para el fundo dominante; (2) falta de uso por diez años; (3) en las personales, la muerte del titular cuando no se fijó plazo; (4) en persona jurídica, vencimiento del plazo máximo de cincuenta años; (5) vencimiento del plazo fijado en el instrumento constitutivo. \\

\textbf{¿Por cuánto tiempo se mantiene la servidumbre de fundo encerrado o sin agua?} &
Mientras se mantengan esas condiciones de encerramiento o la falta de agua. Si desaparece la utilidad (por ejemplo, por la construcción de un camino lateral), desaparece el interés del orden público y la servidumbre carece de objeto. \\

\textbf{¿Qué efectos produce la extinción?} &
Se extinguen todos los derechos personales o reales vinculados con la servidumbre. El dominio del titular del fundo sirviente se restablece como dominio pleno (perfecto), porque desaparece la carga que lo desmembraba. \\

\midrule

\multicolumn{2}{!{\vrule width 0.4pt}>{\raggedright\arraybackslash}X!{\vrule width 0.4pt}}{%
\textbf{Resumen:} La servidumbre se extingue por desaparición de la utilidad, falta de uso por diez años, muerte del titular personal sin plazo fijado, vencimiento del plazo máximo de cincuenta años en persona jurídica o vencimiento del plazo pactado. Al extinguirse, el dominio del fundo sirviente recupera su plenitud.
} \\

\bottomrule
\end{tabularx}
\end{table}

\section{Cuadro Cornell de síntesis general}

\begin{table}[H]
\centering
\begin{tabularx}{\textwidth}{
    !{\vrule width 0.4pt}>{\raggedright\arraybackslash}p{0.26\textwidth}
    !{\vrule width 0.4pt}>{\raggedright\arraybackslash}X
    !{\vrule width 0.4pt}
}
\toprule
\textbf{Preguntas / palabras clave} &
\textbf{Notas / respuestas} \\
\midrule

\textbf{¿Qué vincula siempre la servidumbre y a quién beneficia o grava cada fundo?} &
Siempre vincula dos fundos: el dominante recibe la utilidad y el sirviente la presta, desmembrando su dominio y convirtiéndolo en imperfecto. \\

\textbf{¿Cómo se clasifican las servidumbres?} &
En positivas (soportar) y negativas (abstenerse); en reales (atienden al fundo, se transmiten con el inmueble) y personales (atienden al titular, se extinguen al transmitirse); y, según su origen, en forzosas o legales (orden público) y convencionales (acuerdo de partes). \\

\textbf{¿Qué distingue a la servidumbre forzosa de la convencional?} &
La forzosa se impone por orden público cuando el fundo carece de salida a la vía pública o de agua, y el sirviente solo reclama la indemnización. La convencional nace del acuerdo de las partes, que pactan precio, camino y plazo. \\

\midrule

\multicolumn{2}{!{\vrule width 0.4pt}>{\raggedright\arraybackslash}X!{\vrule width 0.4pt}}{%
\textbf{Resumen:} La servidumbre es un derecho real de disfrute sobre cosa ajena entre dos fundos, que se constituye por escritura pública, se clasifica según la utilidad, el criterio de constitución y el origen, es inembargable e inseparable del fundo (Art. 2178), y se extingue por las causales del Art. 2182, tras lo cual el dominio sirviente recupera su plenitud.
} \\

\bottomrule
\end{tabularx}
\end{table}

\end{document}
